% !Mode:: "TeX:UTF-8" 
\setlength{\subfigbottomskip}{0pt}
\CTEXoptions[bibname={主要参考文献}]
\CTEXsetup[name={,},number={}]{chapter}
\captionsetup{labelsep=space,font=small,justification=centering}
\arraycolsep=1.7pt
\graphicspath{{figures/}}
\renewcommand{\subcapsize}{\zihao{5}}
\renewcommand{\thesubfigure}{\alph{subfigure})}
\setcounter{secnumdepth}{4}
\newcommand{\pozhehao}{\raisebox{0.1em}{------}}
\titleformat{\chapter}{\center\zihao{-2}\heiti}{\chaptertitlename}{0.5em}{}
\titlespacing{\chapter}{0pt}{-4.5mm}{8mm}
\titleformat{\section}{\zihao{-3}\heiti}{\thesection}{0.5em}{}
\titlespacing{\section}{0pt}{4.5mm}{4.5mm}
\titleformat{\subsection}{\zihao{4}\heiti}{\thesubsection}{0.5em}{}
\titlespacing{\subsection}{0pt}{4mm}{4mm}
\titleformat{\subsubsection}{\zihao{-4}\heiti}{\thesubsubsection}{0.5em}{}
\titlespacing{\subsubsection}{0pt}{0pt}{0pt}
\makeatletter
\renewcommand\thesection{\@arabic \c@section} % 前面不带 thechapter
\makeatother

\theoremstyle{plain}
\theorembodyfont{\songti\rmfamily}
\theoremheaderfont{\heiti\rmfamily}
\newtheorem{definition}{\heiti 定义}
\newtheorem{example}{\heiti 例}
\newtheorem{algo}{\heiti 算法}
\newtheorem{theorem}{\heiti 定理}
\newtheorem{axiom}{\heiti 公理}
\newtheorem{proposition}{\heiti 命题}
\newtheorem{lemma}{\heiti 引理}
\newtheorem{corollary}{\heiti 推论}
\newtheorem{remark}{\heiti 注解}
\newenvironment{proof}{\noindent{\heiti 证明：}}{\hfill $ \square $ \vskip 4mm}
\theoremsymbol{$\square$}

% 定义页眉和页脚 使用fancyhdr 宏包
\newcommand{\makeheadrule}{
\rule[7pt]{\textwidth}{0.75pt} \\[-23pt]
\rule{\textwidth}{2.25pt}}
\renewcommand{\headrule}{
    {\if@fancyplain\let\headrulewidth\plainheadrulewidth\fi
     \makeheadrule}}
\makeatother
		
\pagestyle{fancyplain}
\renewcommand{\chaptermark}[1]{\relax}
\renewcommand{\sectionmark}[1]{\markright{#1}}
\fancyhf{}
\fancyhead[C]{\songti \zihao{-5}  哈尔滨工业大学本科毕业设计开题报告}
\fancyfoot[C]{\zihao{-5} -~\thepage~-}
\renewcommand\bibsection{\section*{\centerline{\bibname}}
\markboth{ 哈尔滨工业大学本科毕业设计开题报告}{\bibname}}


\renewcommand{\CJKglue}{\hskip 0.56pt plus 0.08\baselineskip} %加大字间距，使每行33个字
\def\defaultfont{\renewcommand{\baselinestretch}{1.62}\normalsize\selectfont}
% 调整罗列环境的布局
\setitemize{leftmargin=3em,itemsep=0em,partopsep=0em,parsep=0em,topsep=-0em}
\setenumerate{leftmargin=3em,itemsep=0em,partopsep=0em,parsep=0em,topsep=0em}
\renewcommand{\theequation}{\arabic{equation}}
\renewcommand{\thetable}{\arabic{table}}
\renewcommand{\thefigure}{\arabic{figure}}

\makeatletter
\renewcommand{\p@subfigure}{\thefigure~}
\makeatother

\newcommand{\citeup}[1]{\textsuperscript{\cite{#1}}} % for WinEdt users

% 封面、摘要、版权、致谢格式定义
\makeatletter
\def\title#1{\def\@title{#1}}\def\@title{}
\def\titlesec#1{\def\@titlesec{& \rule[-4pt]{200pt}{1pt}\hspace{-326pt}\centerline{\textbf{#1}}}}\def\@titlesec{}
\def\affil#1{\def\@affil{#1}}\def\@affil{}
\def\subject#1{\def\@subject{#1}}\def\@subject{}
\def\author#1{\def\@author{#1}}\def\@author{}
\def\bdate#1{\def\@bdate{#1}}\def\@bdate{}
\def\supervisor#1{\def\@supervisor{#1}}\def\@supervisor{}
\def\date#1{\def\@date{#1}}\def\@date{}
\def\stuno#1{\def\@stuno{#1}}\def\@stuno{}
% %%%%%%%%%%%%%%%%%%%                定义封面                     %%%%%%%%%%%%%%%%%%%%%%%%%
\def\makecover{
    \thispagestyle{empty}
    \zihao{-2}\vspace*{10mm}
		\renewcommand{\CJKglue}{\hskip 2pt plus 0.08\baselineskip}  
\begin{figure}[htbp]
\centering
\includegraphics[scale=0.16]{HIT.eps}
\end{figure}		
\vspace{10mm}
\centerline{\zihao{-1}\songti\textbf{毕业设计（论文）开题报告}}
\renewcommand{\CJKglue}{\hskip 0pt plus 0.08\baselineskip}
\vspace{51.5pt}
\zihao{-2}
\begin{center}\songti\textbf{题~目：\@title}\end{center}
\vspace{120pt}
    \zihao{3}
    \hspace*{68pt}{\songti
	\renewcommand{\arraystretch}{1.8}
    \begin{tabular}{l@{}l}
    \textbf{专\hfill 业}     & \rule[-4pt]{200pt}{1pt}\hspace{-326pt}\centerline{\textbf\@subject}\\
    \textbf{学\hfill 生}      & \rule[-4pt]{200pt}{1pt}\hspace{-326pt}\centerline{\textbf\@author}\\
    \textbf{学\hfill 号}  & \rule[-4pt]{200pt}{1pt}\hspace{-326pt}\centerline{\textbf\@stuno}\\
    \textbf{指\hfill 导\hfill 教\hfill 师}     & \rule[-4pt]{200pt}{1pt}\hspace{-326pt}\centerline{\textbf\@supervisor}\\
    \textbf{日\hfill 期} & \rule[-4pt]{200pt}{1pt}\hspace{-326pt}\centerline{\textbf\@date}\\
    \end{tabular}
		\renewcommand{\arraystretch}{1}}
	\vfill
	\vfill
    \centerline{\songti\textbf{\lishu \zihao{-2}哈尔滨工业大学教务处制}}

%%定义内封
\newpage

%%%%%%%%%%%%%%%%%%%%%%%%%%%%%%%%%%%%%%%%%%%%%%%%%%%%%%%%%%%%%%%%%
\thispagestyle{empty}
\zihao{5}\vspace*{2em}
\begin{center}
  \heiti\zihao{3}说\hspace{3em}明
\end{center}
\vspace*{40pt}
	\renewcommand{\arraystretch}{1.25}
    {\heiti\zihao{4}
    \hangindent=2em
	\noindent {\zihao{-3}\textbf{ 一、开题报告主要内容}}
    \begin{enumerate}[leftmargin=36pt]
	\item 课题来源及研究的目的和意义\\ {\songti \zihao{-4}（正文  宋体小4号字，行距1.25倍，段前0行，段后0行）}
	\item 国内外在该方向的研究现状及分析
	\item 主要研究内容
	\item 研究方案
	\item 进度安排，预期达到的目标
	\item 课题已具备和所需的条件、经费
	\item 研究过程中可能遇到的困难和问题，解决的措施
	\item 主要参考文献
    \end{enumerate}
    \noindent  {\zihao{-3} \textbf{  二、对开题报告的要求}}
	\begin{enumerate}[leftmargin=36pt]
	\item 开题报告的字数应在3000字以上。
	\item 参考文献的要求：		
 \end{enumerate}  
\begin{enumerate}
\renewcommand{\labelenumi}{(\theenumi)} 
\songti \zihao{-4}
\item 理工类论文的参考文献一般为10-15篇，其中学术期刊类文献不少于7篇，外文文献不少于3篇（特殊专业可酌情确定明确要求，并报教务处备案）；文科、管理类论文，参考文献一般为15-20篇，其中学术期刊类文献不少于12篇，外文文献不少于3篇。近五年的文献数不应少于总数的1/3，应有近两年的参考文献。教材、产品说明书、国家标准、未公开发表的研究报告不宜作为参考资料。
\item 参考文献按在开题报告中出现的次序列出。
\item 参考文献书写顺序：序号 作者.文章名.学术刊物名.年，卷（期）：引用起止页。
 \end{enumerate}  
}  
      
%%%%%%%%%%%%%%%%%%%%%%%%%%%%%%%%%%%%%%%%%%%%%%%%%%%%%%%%%%%%%%%%%%%%%%%%%%%%%%%%%%%%%%%
	\renewcommand{\arraystretch}{1}
    \clearpage
}

\makeatother